\section{The benchmark: what the public release contains}
\label{sec:benchmark}

Everything below was established from pinned public artifacts and pinned public
evaluator source. No private attachment, gold SQL, hidden knowledge annotation,
or hidden test case was read. ``Observed'' means reproduced locally; ``official
claim'' means the maintainers state it and this project has not independently
established it.

\paragraph{Pinned sources.}
The dataset is LiveSQLBench Large-v1 at Hugging Face revision
\hash{a418e108d5cbb4cf9b783a928eff5e924ad2460d}, with
\texttt{livesqlbench\_large\_v1\_data.jsonl} verified at SHA-256
\hash{f0e12218cb46f5b6e019908740a0b3303a1f8d1136c661545ad6dd1b4b5444f6}. The
evaluator, database build, and reference agent were read at repository commit
\hash{e15cd221267e06fabfaf6a3d4a69308280ce9a7c}. The dataset card describes the
benchmark as continuously evolving, so the name ``Large-v1'' does not identify
immutable bytes and every experiment record stores both revision and content
hash.

\paragraph{Eligibility.}
The public JSONL holds 480 records with 480 unique \texttt{instance\_id} values.
Filtering on the public \texttt{category} field alone yields exactly 332
\texttt{Query} records and 148 \texttt{Management} records. \texttt{Query} is the
SELECT-only analytical class; \texttt{Management} covers database-management and
CRUD tasks, outside the scope of a semantic-layer comparison, and is excluded.
All 18 databases contribute, with per-database eligible counts from 10 to 24.

Public records carry \texttt{instance\_id}, \texttt{selected\_database},
\texttt{query}, \texttt{normal\_query}, \texttt{category}, \texttt{high\_level},
\texttt{conditions}, \texttt{preprocess\_sql}, and \texttt{clean\_up\_sqls}. The
protected fields \texttt{sol\_sql}, \texttt{external\_knowledge}, and
\texttt{test\_cases} appear as empty arrays in all 480 records; they are not
omitted. The \texttt{difficulty\_tier} field documented on the card does not exist
at this revision, so difficulty is used neither for splitting nor for subgroup
analysis. Two naming inconsistencies require adapters, not assumptions:
the released record uses \texttt{clean\_up\_sqls} while both card and pinned
evaluator use \texttt{clean\_up\_sql}. All 332 eligible records have empty
\texttt{preprocess\_sql} and \texttt{clean\_up\_sqls} arrays here, so the
mismatch does not change current inputs, but the adapter handles the contract
explicitly instead of relying on that emptiness persisting.

Public metadata available for stratification: \texttt{high\_level} is 181 true
and 151 false; \texttt{conditions.order} 194 true and 138 false;
\texttt{conditions.distinct} 17 true and 315 false; and
\texttt{conditions.decimal} takes values $-1$ (158), 0 (8), 1 (6), 2 (110), 3
(24), 4 (22), 5 (1), 6 (2), and 8 (1). Database is the dominant stratum. Several
joint \texttt{conditions} cells are too sparse to cross-stratify without building
strata small enough to determine individual questions, so those marginals are
audited and reported after allocation rather than used to construct strata.

\paragraph{Scale and schema metadata.}
The official documentation reports 18 PostgreSQL databases, 971 tables, 17,749
columns, roughly 2,056 rows per table on average, and 746.33\,MB total. Counting
\texttt{CREATE TABLE} statements across pinned public schema files independently
reproduces 971 tables, and the column-meaning maps contain exactly 17,749 keys.
Two properties shape the transformation work. First, 212 of those values are
structured, not plain strings: for JSON and JSONB columns they carry a
column-level description plus a \texttt{fields\_meaning} map describing nested
properties, so a transformation that flattens only relational columns discards
the definitions that JSON operations need. Second, the schema files are not clean
DDL; they interleave \texttt{CREATE TABLE} statements with three sample rows per
table and retain original quoting, casing, and key declarations, so preparation
must parse DDL separately from examples while preserving both as inputs with
different provenance.

\paragraph{The knowledge base is a dependency graph.}
Each database ships a \texttt{*\_kb.jsonl} whose records carry an integer
\texttt{id} unique within the database, a \texttt{knowledge} concept name, a
\texttt{description}, a \texttt{definition} giving the rule or formula in natural
language, a \texttt{type}, and \texttt{children\_knowledge}. Despite the name,
\texttt{children\_knowledge} lists the identifiers the entry \emph{depends on},
with $-1$ documented as no dependency. Table~\ref{tab:hkb} gives the observed
graph statistics; both $-1$ and the undocumented empty list are normalized to
zero outgoing edges.

\begin{table}[htbp]
\centering
\caption{Observed hierarchical knowledge base structure across all 18 databases.}
\label{tab:hkb}
\begin{tabular}{lr}
\toprule
Property & Count \\
\midrule
HKB entries (\texttt{calculation} 430, \texttt{domain} 462, \texttt{value\_illustration} 198) & 1,090 \\
Entries using the documented no-dependency sentinel $-1$ & 509 \\
Entries using the undocumented empty list & 21 \\
Entries with one or more direct dependencies & 560 \\
Direct dependency edges & 945 \\
Maximum direct dependencies on one entry & 5 \\
Edges whose target itself has dependencies & 344 \\
Databases containing multi-hop dependencies & 18 of 18 \\
Maximum dependency-chain length & 6 edges \\
Duplicate identifiers / dangling refs / self-edges / cycles & 0 / 0 / 0 / 0 \\
\bottomrule
\end{tabular}
\end{table}

This is what makes the benchmark suitable for the question. An HKB entry cannot
be treated as an isolated blob of prompt text: 344 edges lead to another derived
entry and multi-hop resolution is required in every database. A transformation
into a semantic layer must preserve a stable source key, the dependency edges,
the original type, the source text, and the compiled objects implementing each
entry, and it can be evaluated on whether that composition survives.

Two documentation discrepancies are recorded rather than resolved by inference.
The general card describes two HKB formats, structured JSON and an unstructured
document, while this snapshot contains only \texttt{*\_kb.jsonl} and calls itself
HKB-JSON. The card also advertises business-rule drift with versioned
definitions, while observed records carry no version or validity field. Whether
drift is encoded only in text, deferred to a later release, or present in
protected data is unresolved, and will not be inferred from gold.

\paragraph{Official evaluator behavior.}
For \texttt{category == "Query"} the runner discards any customized test-case
list and installs one default Soft EX result-equivalence test. At the pinned
revision that path removes block and newline-terminated line comments; removes
standalone \texttt{DISTINCT} while attempting to preserve \texttt{DISTINCT ON};
replaces each \texttt{ROUND(expression, precision)} with its first argument
including nested calls; executes both statements under a 60-second PostgreSQL
statement timeout; fetches at most 10,000 rows; normalizes top-level
\texttt{date} and \texttt{datetime} values to \texttt{YYYY-MM-DD}, recursively
rounds floats and decimals to two places, and serializes lists and dictionaries
with sorted keys; requires both normalized results to be non-empty; and compares
normalized sequences exactly when \texttt{conditions.order} is true and Python
sets of rows otherwise.

Consequences that are load-bearing for interpretation follow directly. Unordered
comparison discards duplicate multiplicity. Two empty result sets fail rather
than match. Date-time normalization discards time of day. The default path
hard-codes two decimal places and never consults \texttt{conditions.decimal}, and
rewrites \texttt{DISTINCT} regardless of \texttt{conditions.distinct}. The scorer
executes the prediction once to detect errors and again inside the comparison,
and for a list of statements retains only the last statement's result. These come
from the pinned implementation, not the README description, and any sealed
reimplementation must reproduce or explicitly test each.

\paragraph{Reference agent and the knowledge-access boundary.}
The repository also ships LiveSQLBench-Agent, an orchestrator with eight tools, a
30-step budget, SELECT-only exploration, one final SQL submission, and per-task
database copies. It is a useful harness skeleton because schema inspection,
column meanings, knowledge discovery, execution, and submission are separate
operations. It is not usable unchanged: its documented configuration accepts only
\texttt{lite} or \texttt{full} at this revision, and its knowledge endpoint omits
\texttt{type} and \texttt{children\_knowledge} from what the agent sees, erasing
precisely the dependency structure this study evaluates. It is nonetheless the
best public evidence for the intended access boundary. Its tools let a model list
database-level knowledge names, fetch a named definition, or fetch all
definitions, and its instruction tells the model to discover relevant knowledge;
no tool receives hidden per-question \texttt{external\_knowledge} identifiers as
a runtime oracle. This protocol follows the same distinction for C2 through C4.

\paragraph{Provisioning and parity.}
The official from-scratch path downloads the Large-v1 dumps, builds a PostgreSQL
18 image, imports 18 template databases, and clones from them; the prebuilt
compose file declares no Large-v1 service, so the from-source build is the
verified official path. The scoring authority is the official local snapshot even
if the evaluated product must query a managed mirror. Because a mismatch between
generation and scoring copies can manufacture both false gains and false
failures, schema fingerprints and deterministic data canaries are compared across
copies before any run, and the parity check reports pass, observed disagreement,
and could-not-run as three distinct states, so missing access can never be
recorded as a passed or failed benchmark question.
